\begin{tikzpicture}[
    scale=1.2,
    % Define a style that places an arrowhead exactly halfway along the path
    mid arrow/.style={
        postaction={
            decorate,
            decoration={
                markings,
                mark=at position 0.67 with {\arrow{Stealth[scale=2.2, fill=black]}}
            }
        }
    }
]

    % Securely define the green/olive/white mix once before rendering anything
    \colorlet{MyLightOlive}{green!50!olive!62!white}

    % Central point of the grid
    \coordinate (Center) at (0.5, 0.5);

    % --- PASS 1: Draw the background grid (Bottom Layer) ---

    % Level 1: Draw ALL 49 squares as GREEN by default
    \foreach \x in {-3,...,3} {
        \foreach \y in {-3,...,3} {
            \filldraw[fill=MyLightOlive, draw=black] (\x,\y) rectangle ++(1,1);
        }
    }
    
    % Level 2: Draw the inner 3x3 (9 squares) as WHITE on top of the green
    \foreach \x in {-1,0,1} {
        \foreach \y in {-1,0,1} {
            \filldraw[fill=white, draw=black] (\x,\y) rectangle ++(1,1);
        }
    }
    
    % Level 3: Draw the single center square as GRAY on top of the white
    \filldraw[fill=lightgray, draw=black] (0,0) rectangle ++(1,1);

    % --- PASS 2: Draw the circles (Middle Layer) ---
    \foreach \x in {-3,...,3} {
        \foreach \y in {-3,...,3} {
            
            \coordinate (SquareCenter) at (\x+0.5, \y+0.5);
            \pgfmathsetmacro{\CircleRadius}{sqrt(2)/2}
            
            \ifthenelse{\x=0 \AND \y=0}{
                % Draw solid circle on the center gray square
                \draw[black, thin] (SquareCenter) circle (\CircleRadius);
            }{
                \ifthenelse{\x>1 \OR \x<-1 \OR \y>1 \OR \y<-1}{
                     % Draw solid circle on the outer green squares
                    \draw[black, thin] (SquareCenter) circle (\CircleRadius);
                }{}
            }
        }
    }

    % --- PASS 3: Draw the spindles, arrows, and dots (Top Layer) ---
    \foreach \x in {-3,...,3} {
        \foreach \y in {-3,...,3} {
            
            \coordinate (SquareCenter) at (\x+0.5, \y+0.5);
            
            % Only draw dots and lines for the green squares (outside the 3x3 ring)
            \ifthenelse{\x>1 \OR \x<-1 \OR \y>1 \OR \y<-1}{
                
                % We draw the dot  
                \filldraw[black] (SquareCenter) circle (1.5pt);
                
                % Draw the TRUE spindle shape for ALL green squares, making it fully black
                \path (SquareCenter) -- (Center) coordinate[pos=0.5] (Mid);
                \fill[black, draw=black, line width=0.1pt] 
                    (SquareCenter) 
                    -- ($ (Mid)!0.8pt!90:(Center) $) 
                    -- (Center) 
                    -- ($ (Mid)!0.8pt!-90:(Center) $) 
                    -- cycle;
                
                % Check if the current square is one of the 8 specific outer squares
                \pgfmathsetmacro{\isCorner}{abs(\x)==3 && abs(\y)==3 ? 1 : 0}
                \pgfmathsetmacro{\isEdgeCenter}{(abs(\x)==3 && \y==0) || (\x==0 && abs(\y)==3) ? 1 : 0}
                
                \pgfmathparse{\isCorner + \isEdgeCenter > 0 ? 1 : 0}
                \ifnum\pgfmathresult=1
                    % Draw an invisible path using the 'mid arrow' style to place the arrowhead
                    \path[mid arrow] (SquareCenter) -- (Center);
                \fi
            }{}
        }
    }
    
    % Re-draw the central dot on the very top so it isn't covered by the lines
    \filldraw[black] (Center) circle (2.5pt);

    % --- PASS 4: Custom Legend ---
    \begin{scope}[shift={(0, -6)}]
        \pgfmathsetmacro{\CircleRadius}{sqrt(2)/2}
        
        % Legend Item 1: Multipoles (Green cell, circle, dot) using pre-defined safe color
        \filldraw[fill=MyLightOlive, draw=black] (-2.5,-0.5) rectangle ++(1,1);
        \draw[black, thin] (-2,0) circle (\CircleRadius);
        \filldraw[black] (-2,0) circle (2pt);
        \node[right, align=left] at (-1.2, 0) {Multipole\\moments};

        % Legend Item 2: Local Moment (Gray cell, circle, dot)
        \filldraw[fill=lightgray, draw=black] (1.0,-0.5) rectangle ++(1,1);
        \draw[black, thin] (1.5,0) circle (\CircleRadius);
        \filldraw[black] (1.5,0) circle (3pt);
        \node[right, align=left] at (2.3, 0) {Local\\moments};

    \end{scope}

\end{tikzpicture}